\section{9.1 퍼셉트론의 한계, 순전파, 역전파}
\begin{frame}{퍼셉트론: ``AI의 시작''이자 ``AI의 실망''}
  \begin{columns}[T]
    \begin{column}{0.55\textwidth}
      \begin{itemize}
        \item \kb{1958, 로젠블라트(코넬)}: \kb{퍼셉트론} 발표 ---
          ``컴퓨터가 배운다'', ``우주 탐사 1세대 로봇''
        \item Week 2의 로지스틱회귀와 거의 동일, 시그모이드 대신 \kb{스텝함수}
          \[\hat{y} = \operatorname{sgn}(w^\top x + b)\]
        \item 실제로 (당시 기준) 조종사 얼굴 인식 실험까지 수행
      \end{itemize}
    \end{column}
    \begin{column}{0.43\textwidth}
      \centering
      \includegraphics[width=\linewidth]{ch09_1_sgn_vs_sigmoid.png}
      {\footnotesize sgn(파랑, $z=0$에서 도약) vs 시그모이드(주황, 매끄러운 전이)}
    \end{column}
  \end{columns}
  \vskip0.3em
  {\small \kb{퍼셉트론 학습 규칙} (경사하강법의 선조): 틀린 예제에만 오차 방향으로 갱신,
  $w \leftarrow w + y\,x$ — 선형분리 가능 데이터에서 \textbf{반드시 수렴}(수렴 정리)}
\end{frame}
